\documentclass[10pt,landscape,a4paper]{article}

% --- PAQUETS ET CONFIGURATION ---
\usepackage[utf8]{inputenc}
\usepackage[T1]{fontenc}
\usepackage[french]{babel}
\usepackage[margin=1cm]{geometry} % Marges fines
\usepackage{multicol} % Pour les colonnes
\usepackage{xcolor}
\usepackage{enumitem}
\usepackage{listings} % Pour le code C
% \usepackage{fontawesome5} % Pour les petites icônes (optionnel, sinon retirer)
\usepackage[most]{tcolorbox} % Pour les jolies boîtes

% --- COULEURS PERSONNALISÉES ---
\definecolor{draculaBg}{RGB}{40, 42, 54}
\definecolor{draculaFg}{RGB}{248, 248, 242}
\definecolor{codeKeyword}{RGB}{255, 121, 198}
\definecolor{codeString}{RGB}{241, 250, 140}
\definecolor{codeComment}{RGB}{98, 114, 164}
\definecolor{boxBlue}{RGB}{50, 150, 250}
\definecolor{boxGreen}{RGB}{80, 200, 120}
\definecolor{boxRed}{RGB}{220, 50, 50}
\definecolor{boxPurple}{RGB}{150, 50, 200}
\definecolor{boxOrange}{RGB}{255, 140, 0}

% --- CONFIGURATION DU CODE C ---
\lstset{
    language=C,
    basicstyle=\ttfamily\scriptsize,
    keywordstyle=\color{codeKeyword}\bfseries,
    stringstyle=\color{codeString},
    commentstyle=\color{codeComment}\itshape,
    backgroundcolor=\color{gray!10},
    frame=none,
    breaklines=true,
    tabsize=2,
    showstringspaces=false
}

% --- STYLE DES BOÎTES ---
\newtcolorbox{mybox}[2][]{
    enhanced,
    colback=white,
    colframe=#2,
    coltitle=white,
    fonttitle=\bfseries\sffamily,
    title=#1,
    arc=3mm,
    boxrule=0.5mm,
    drop shadow,
    left=1mm, right=1mm, top=1mm, bottom=1mm,
    before skip=2mm, after skip=2mm
}

% --- DÉBUT DU DOCUMENT ---
\begin{document}

% Titre stylé
\begin{center}
    {\Huge \sffamily \textbf{Survival Kit : Langage C \& Chaînes}} \\
    {\small \textit{L'antisèche ultime pour les partiels - Gestion de la mémoire et manipulation de texte}}
\end{center}

\begin{multicols*}{3}

% ----------------------------------------------------
% SECTION 1: STRING.H - LES BASIQUES
% ----------------------------------------------------
\begin{mybox}[BASIQUES <string.h>]{boxBlue}
    \textit{Opérations classiques. Attention au \texttt{\textbackslash0} !}
    \begin{itemize}[leftmargin=*, noitemsep]
        \item \textbf{\texttt{size\_t strlen(s)}} : Longueur sans \texttt{\textbackslash0}.
        \item \textbf{\texttt{strcpy(dest, src)}} : Copie \texttt{src} dans \texttt{dest}.
            \textit{\textcolor{red}{Danger:}} Buffer overflow si dest trop petit.
        \item \textbf{\texttt{strncpy(dest, src, n)}} : Copie max \texttt{n} char.
            \textit{\textcolor{red}{Piège:}} Ne met pas de \texttt{\textbackslash0} si coupé !
        \item \textbf{\texttt{strcat(dest, src)}} : Ajoute à la fin.
        \item \textbf{\texttt{strncat(dest, src, n)}} : Ajout sécurisé.
        \item \textbf{\texttt{strcmp(s1, s2)}} : Compare.
            \begin{itemize}
                \item \texttt{0} : Égales.
                \item \texttt{<0} : s1 avant s2.
                \item \texttt{>0} : s1 après s2.
            \end{itemize}
        \item \textbf{\texttt{strdup(s)}} : Malloc + Copie + Return.
            \textit{\textcolor{boxRed}{Impératif:}} Faire un \texttt{free()} après.
    \end{itemize}
\end{mybox}

% ----------------------------------------------------
% SECTION 2: RECHERCHE & ANALYSE
% ----------------------------------------------------
\begin{mybox}[RECHERCHE <string.h>]{boxBlue}
    \begin{itemize}[leftmargin=*, noitemsep]
        \item \textbf{\texttt{strchr(s, c)}} : Pointe vers la 1ère occurence de \texttt{c}. Return \texttt{NULL} si absent.
        \item \textbf{\texttt{strrchr(s, c)}} : Pointe vers la \textbf{dernière} occurence.
        \item \textbf{\texttt{strstr(haystack, needle)}} : Cherche une sous-chaîne dans une chaîne.
        \item \textbf{\texttt{strtok(str, delim)}} : Découpe la chaîne (split).
            \textit{Note:} Destructif (remplace sép. par \texttt{\textbackslash0}). Utiliser \texttt{NULL} pour continuer.
            \begin{lstlisting}
char *token = strtok(str, " ");
while (token) {
    printf("%s\n", token);
    token = strtok(NULL, " ");
}
            \end{lstlisting}
    \end{itemize}
\end{mybox}

% ----------------------------------------------------
% SECTION 3: INSPECTION (CTYPE.H)
% ----------------------------------------------------
\begin{mybox}[INSPECTION <ctype.h>]{boxGreen}
    \textit{Test sur un seul caractère (int c).}
    \begin{itemize}[leftmargin=*, noitemsep]
        \item \textbf{\texttt{isalpha(c)}} : A-Z ou a-z.
        \item \textbf{\texttt{isdigit(c)}} : 0-9.
        \item \textbf{\texttt{isalnum(c)}} : Lettre OU Chiffre.
        \item \textbf{\texttt{isspace(c)}} : Espace, \texttt{\textbackslash t}, \texttt{\textbackslash n}, \texttt{\textbackslash v}, \texttt{\textbackslash f}, \texttt{\textbackslash r}.
        \item \textbf{\texttt{isprint(c)}} : Caractère imprimable.
        \item \textbf{\texttt{toupper(c)} / \texttt{tolower(c)}} : Conversion casse.
    \end{itemize}
\end{mybox}

% ----------------------------------------------------
% SECTION 4: CONVERSIONS
% ----------------------------------------------------
\begin{mybox}[CONVERSIONS]{boxOrange}
    \textit{Bibliothèques <stdlib.h> et <stdio.h>}
    \begin{itemize}[leftmargin=*, noitemsep]
        \item \textbf{\texttt{atoi(str)}} : String vers Int.
            \textit{Limité:} Retourne 0 si erreur.
        \item \textbf{\texttt{strtol(str, \&end, base)}} : La version Pro.
            Permet de gérer les erreurs et la base (10, 16, 2...).
        \item \textbf{\texttt{sprintf(buf, fmt, ...)}} : "Print" dans une string.
            \textit{Ex:} \texttt{sprintf(buf, "\%d", 42);} $\rightarrow$ Int to Str.
        \item \textbf{\texttt{sscanf(str, fmt, ...)}} : Lire depuis une string.
    \end{itemize}
\end{mybox}

% ----------------------------------------------------
% SECTION 5: MÉMOIRE BRUTE
% ----------------------------------------------------
\begin{mybox}[MÉMOIRE <string.h>]{boxRed}
    \textit{Opérations sur \texttt{void*}. Pas d'arrêt au \texttt{\textbackslash0}.}
    \begin{itemize}[leftmargin=*, noitemsep]
        \item \textbf{\texttt{memset(ptr, val, n)}} : Remplit \texttt{n} octets avec \texttt{val}.
            \textit{Usage:} \texttt{memset(arr, 0, sizeof(arr));}
        \item \textbf{\texttt{memcpy(dst, src, n)}} : Copie rapide.
            \textit{\textcolor{red}{Interdit:}} Chevauchement (overlap).
        \item \textbf{\texttt{memmove(dst, src, n)}} : Copie sûre (gère l'overlap).
        \item \textbf{\texttt{memcmp(p1, p2, n)}} : Compare deux blocs mémoire.
    \end{itemize}
\end{mybox}

% ----------------------------------------------------
% SECTION 6: ALGOS CLASSIQUES
% ----------------------------------------------------
\begin{mybox}[ALGOS À LA CON (EXAM)]{boxPurple}
    \textbf{1. Swap sans variable temp}
    \begin{lstlisting}
void swap(int *a, int *b) {
    *a = *a ^ *b;
    *b = *a ^ *b;
    *a = *a ^ *b;
}
    \end{lstlisting}
    
    \textbf{2. Inverser une chaîne (in-place)}
    \begin{lstlisting}
void strrev(char *str) {
    int i = 0, j = strlen(str) - 1;
    while (i < j) {
        char temp = str[i];
        str[i++] = str[j];
        str[j--] = temp;
    }
}
    \end{lstlisting}
    
    \textbf{3. Est puissance de 2 ?}
    \begin{lstlisting}
int is_pow2(int n) {
    return (n > 0 && !(n & (n - 1)));
}
    \end{lstlisting}
\end{mybox}

% ----------------------------------------------------
% SECTION 7: POINTEURS & PIÈGES
% ----------------------------------------------------
\begin{mybox}[MEMO POINTEURS]{boxRed}
    \begin{itemize}[leftmargin=*, noitemsep]
        \item \textbf{\texttt{char *s = "Bonjour";}} :
            Stocké en lecture seule (RO). \texttt{s[0] = 'b'} $\rightarrow$ \textbf{Segfault}.
        \item \textbf{\texttt{char s[] = "Bonjour";}} :
            Tableau modifiable sur la stack. OK.
        \item \textbf{\texttt{malloc(n)}} :
            Toujours vérifier si retour != NULL.
        \item \textbf{\texttt{free(ptr)}} :
            Ne jamais free deux fois (Double Free). Mettre \texttt{ptr = NULL} après.
    \end{itemize}
\end{mybox}

% ----------------------------------------------------
% SECTION 8: FONCTIONS BONUS
% ----------------------------------------------------
\begin{mybox}[BONUS (POUR BRILLER)]{boxOrange}
    \begin{itemize}[leftmargin=*, noitemsep]
        \item \textbf{\texttt{strcspn(s, reject)}} : Longueur du segment initial qui NE contient PAS les chars de \textit{reject}.
            \textit{Astuce:} Pour enlever le \texttt{\textbackslash n} de \texttt{fgets} :
            \texttt{s[strcspn(s, "\textbackslash n")] = 0;}
        \item \textbf{\texttt{strspn(s, accept)}} : Longueur du segment initial contenant QUE des chars de \textit{accept}.
        \item \textbf{\texttt{perror("msg")}} : Affiche l'erreur système actuelle sur stderr.
    \end{itemize}
\end{mybox}

\end{multicols*}
\end{document}